\documentclass[11pt,a4paper]{article}
\usepackage[margin=1in]{geometry}
\usepackage[T1]{fontenc}
\usepackage{lmodern}
\usepackage{microtype}
\usepackage{xcolor}
\usepackage{booktabs}
\usepackage{amsmath}
\usepackage{hyperref}
\usepackage{listings}
\usepackage{enumitem}
\usepackage{fancyhdr}
\definecolor{codebg}{RGB}{245,247,250}
\lstset{basicstyle=\ttfamily\small,backgroundcolor=\color{codebg},frame=single,breaklines=true,columns=fullflexible,keepspaces=true}
\hypersetup{colorlinks=true,linkcolor=blue,urlcolor=blue}
\pagestyle{fancy}
\fancyhf{}
\lhead{Unified JEPA}
\rhead{Original Architecture}
\cfoot{\thepage}
\title{Unified JEPA --- Original Present Architecture}
\author{Metasurface Inverse Design}
\date{September 2026}
\begin{document}
\maketitle
\tableofcontents
\newpage

\section{Scope}
This document describes the original repaired architecture identified as
\texttt{unified\_occ\_param\_spectrum\_jepa\_v1}. It intentionally excludes the
later experimental additions: the direct goal residual, real-vs-null ranking
loss, and detached comparison branches.

The task is spectrum-conditioned inverse design:
\begin{lstlisting}
incomplete geometry + scalar observations + requested spectrum
                           |
                           v
                    completed geometry
                           |
                           v
                  simulated target spectrum
\end{lstlisting}
The model must reconstruct valid geometry and change that geometry in response
to the requested target spectrum.

\section{Semantic data representation}
The released MetaDiT geometry has shape \texttt{[3,64,64]} and uses:
\begin{lstlisting}
channel 0 = occupancy * (r_atom / 5)
channel 1 = occupancy * h_atom
channel 2 = occupancy * (l_lattice / 3)
\end{lstlisting}
The unified model factorizes this into:
\begin{lstlisting}
occupancy: [1,64,64]
scalars:   [l_lattice, h_atom, r_atom]
\end{lstlisting}
The scalar order is fixed everywhere. Approximate valid ranges are:
\begin{lstlisting}
l_lattice: [2.5, 3.0]
h_atom:    [0.5, 1.0]
r_atom:    [3.5, 5.0]
\end{lstlisting}
Every scalar also has a known/unknown flag. Unknown values are zeroed before
the scalar encoder. Known values are hard-retained during final assembly.

\section{Occupancy masking}
The \texttt{64x64} occupancy image is represented on a \texttt{16x16} token
grid with patch size 4. Training masks one to four axis-aligned rectangular
blocks using ratios 25\%, 50\%, 75\%, and 100\%.
\begin{lstlisting}
1 = visible token
0 = masked token
\end{lstlisting}
Masked positions receive learned mask tokens and positional information.
Visible occupancy pixels are hard-retained during geometry assembly. The
critical evaluation regime is:
\begin{lstlisting}
100% occupancy masked + all scalars unknown
\end{lstlisting}

\section{Original end-to-end data flow}
\begin{lstlisting}
occupancy + scalar values/flags
              |
              v
       masked occupancy encoder
              |
       geometry tokens z_x
              |------------------+
              v                  v
       scalar encoder      released spectrum encoder
              |                  |
       scalar summary      local spectrum tokens
              |                  |
              |           c_physics + a_goal
              +----------+-------+
                         |
                         v
                  fusion transformer
                         |
                         v
                    GCLCT predictor
                         |
                         v
                       z_hat
                         |
                         v
                    geometry decoder
                         |
                         v
                   assembled geometry
                         |
                         v
                  frozen EM surrogate
                         |
                         v
                   predicted spectrum
\end{lstlisting}
The EMA target path is separate:
\begin{lstlisting}
complete true geometry -> EMA encoder -> z_y
\end{lstlisting}
The EMA target receives geometry but not the requested spectrum.

\section{Geometry and scalar encoders}
The masked occupancy encoder maps the input into 256 geometry tokens:
\begin{lstlisting}
[B,1,64,64] -> patch embedding -> [B,256,192] = z_x
\end{lstlisting}
The scalar encoder receives:
\begin{lstlisting}
[l_value, l_known,
 h_value, h_known,
 r_value, r_known]
\end{lstlisting}
It produces FiLM parameters for scalar-conditioned processing and one scalar
summary token for the fusion and predictor path.

\section{Released spectrum path}
The released MetaDiT spectrum encoder remains frozen:
\begin{lstlisting}
S [B,2,301] -> released encoder -> A_local [B,301,256]
\end{lstlisting}

\subsection{Global physics condition}
The local spectrum tokens are mean-pooled over frequency and projected:
\begin{lstlisting}
A_local.mean(frequency) -> c_physics [B,384]
\end{lstlisting}
\texttt{c\_physics} FiLM-modulates every GCLCT predictor block.

\subsection{Structured goal tokens}
Sixteen learned queries cross-attend over the 301 local spectrum tokens:
\begin{lstlisting}
learned queries + A_local -> a_goal [B,16,384]
\end{lstlisting}
These tokens enter both the fusion sequence and predictor key/value sequence.

\subsection{Null diagnostic}
The null goal mode replaces both spectrum representations with zeros:
\begin{lstlisting}
c_physics = 0
a_goal = 0
\end{lstlisting}
This tests whether the model uses the requested spectrum or falls back to a
geometry prior. The released spectrum encoder itself was not redesigned.

\section{Fusion and GCLCT predictor}
The fusion transformer concatenates 256 occupancy tokens, 16 goal tokens, and
one scalar token, for 273 tokens total.

The GCLCT predictor receives 256 occupancy completion queries plus one scalar
query, a key/value sequence containing 256 fused occupancy tokens plus 16 goal
tokens, and the global condition \texttt{c\_physics}.

Each block contains:
\begin{lstlisting}
affine-less LayerNorm -> FiLM(c_physics) -> self-attention
affine-less LayerNorm -> FiLM(c_physics) -> cross-attention
affine-less LayerNorm -> FiLM(c_physics) -> MLP
\end{lstlisting}
The output is \texttt{z\_hat [B,257,192]}. The first 256 tokens predict
occupancy latents; the final token predicts the scalar summary.

\section{EMA target and JEPA/VICReg objective}
The complete true occupancy is passed through a frozen EMA copy:
\begin{lstlisting}
complete occupancy -> EMA encoder -> z_y_raw [B,256,192]
\end{lstlisting}
The EMA scalar copy supplies target-side FiLM conditioning. It is not a second
scalar latent target.

The objective owns a shared projector:
\begin{lstlisting}
z_hat -> projector -> p_hat
z_y   -> projector -> p_y
\end{lstlisting}
The representation losses are:
\begin{lstlisting}
L_inv = masked prediction invariance
L_var = variance regularization
L_cov = covariance regularization
\end{lstlisting}
The repaired weights were approximately \texttt{lambda\_inv=25},
\texttt{lambda\_var=25}, and \texttt{lambda\_cov=1}.

The decoder consumes raw \texttt{z\_hat}, while the dominant representation
objective is measured in projected space:
\begin{lstlisting}
raw latent cosine       = -0.06618
projected latent cosine =  0.997437
\end{lstlisting}
This is the projector-absorption risk: projected alignment can look excellent
while the raw decoder latent remains poorly aligned.

\section{Scalar and occupancy decoders}
The scalar-summary token is decoded into the three physical parameters. Each
output is bounded with:
\begin{lstlisting}
prediction = lo + (hi - lo) * sigmoid(raw)
\end{lstlisting}
The 256 occupancy tokens are reshaped into a \texttt{16x16} feature map,
upsampled to \texttt{64x64}, and decoded into occupancy logits. The decoder is
FiLM-conditioned by effective scalars:
\begin{lstlisting}
known scalar   -> true scalar
unknown scalar -> predicted scalar
\end{lstlisting}
The active occupancy reconstruction term is:
\begin{lstlisting}
L_occ = BCEWithLogits(occupancy_logits, true_occupancy)
\end{lstlisting}
It is reported with occupancy IoU, occupancy F1, predicted occupancy fraction,
and true occupancy fraction.

\section{Geometry assembly and physics surrogate}
The decoded outputs are assembled back to the MetaDiT convention:
\begin{lstlisting}
occupancy * (r_atom / 5)
occupancy * h_atom
occupancy * (l_lattice / 3)
\end{lstlisting}
Visible occupancy pixels and known scalar values are retained before the
surrogate.

The assembled geometry is passed through the frozen MetaDiT forward surrogate:
\begin{lstlisting}
geometry [B,3,64,64] -> frozen surrogate -> spectrum [B,2,301]
\end{lstlisting}
The surrogate stays in evaluation mode and receives no parameter gradients.
Gradients with respect to geometry remain enabled. The normalized physics loss
is a SmoothL1 comparison between predicted and requested spectra.

\section{Repaired original objective}
The repaired objective is:
\begin{lstlisting}
L = lambda_inv    * L_inv
  + lambda_var    * L_var
  + lambda_cov    * L_cov
  + lambda_scalar * L_scalar
  + lambda_occ    * L_occ
  + lambda_phys   * L_phys
\end{lstlisting}

The key repairs were:
\begin{itemize}[leftmargin=*]
\item \textbf{Occupancy supervision:} adding \texttt{L\_occ} so the
\texttt{lambda\_phys=0} baseline also trains its occupancy decoder.
\item \textbf{Validation physics:} allowing the real surrogate path to run
during validation even when the model is in \texttt{eval()} mode.
\item \textbf{Scalar validity:} bounding scalar outputs and reporting violations.
\item \textbf{Reproducibility:} recording joint regimes, diagnostics,
validation IDs, checkpoint manifests, and SHA-256 hashes.
\end{itemize}

\section{Repaired original-architecture result}
The repaired original architecture used a selected physics weight near
\texttt{lambda\_phys=0.2}, seed 42, and the hard stratum of 100\% occupancy
masking with all scalars unknown.

\begin{center}
\begin{tabular}{lr}
\toprule
Metric & Value \\
\midrule
Physics real & 0.492732 \\
Physics null & 0.485704 \\
Physics shuffled & 0.518057 \\
Occupancy IoU & 0.690314 \\
Occupancy F1 & 0.815067 \\
Scalar normalized MAE & 0.237963 \\
Scalar out-of-range fraction & 0 \\
Geometry sensitivity & 0.0111547 \\
$c_{\rm physics}$ variation & 0.091998 \\
$a_{\rm goal}$ variation & 0.085290 \\
\bottomrule
\end{tabular}
\end{center}

Lower physics error is better. Geometry reconstruction is non-degenerate,
scalar predictions are valid, and real goals outperform shuffled goals.
However, real goals do not outperform null goals, so target conditioning is
measurable but not yet reliable.

\section{Later experiments excluded from this architecture}
After the original architecture, three separate experiments were tested.

\subsection{Direct goal residual}
The experiment added:
\begin{lstlisting}
z_base + direct spectrum residual = z_final
\end{lstlisting}
This gave spectrum information a shorter path to masked occupancy tokens. The
first version increased goal sensitivity but still had a spectrum-aware base
predictor, so it was not a fully separated base/goal design.

\subsection{Real-vs-null ranking}
A ranking term required real physics error to be lower than null and shuffled
errors. The first version could satisfy the ordering by making null outputs
worse, so it was not accepted as a production architecture.

\subsection{Detached comparison branches}
The next version detached null and shuffled comparison errors so the ranking loss
could not improve merely by damaging comparison outputs. This remains an
experiment and is not part of the original architecture in this document.

\section{Strengths and limitation}
The original design has solid foundations: semantic factorization, explicit
scalar missingness, rectangular block masking, released encoder reuse, an EMA
geometry target, a frozen but differentiable surrogate, hard retention of known
values, bounded scalar outputs, active occupancy reconstruction, curriculum
logging, gradient ownership checks, and recoverable checkpoints.

The main limitation is not raw data factorization. Geometry and spectrum are
already separate semantic inputs. The limitation is that they are mixed before
the main prediction objective is resolved:
\begin{lstlisting}
geometry tokens + goal tokens + scalar token
              -> shared predictor
\end{lstlisting}
At the same time, JEPA targets a geometry-only latent. Therefore the
spectrum-conditioned student is trained toward a target that contains no
spectrum information. This explains the combination of strong projected JEPA
alignment, weak raw decoder-latent alignment, nonzero spectrum sensitivity, and
a weak real-goal advantage over null.

\section{One-page summary}
\begin{lstlisting}
DATA
  MetaDiT geometry -> occupancy + [l_lattice,h_atom,r_atom]
  target spectrum  -> frozen released spectrum encoder

CONTEXT
  block-masked occupancy -> 256 geometry tokens
  scalar values/flags    -> FiLM + scalar summary

GOAL
  spectrum tokens -> c_physics + 16 a_goal tokens

PREDICTOR
  geometry + goal + scalar tokens -> fusion -> GCLCT -> z_hat

TARGET
  complete geometry -> EMA encoder -> z_y

DECODER
  z_hat -> occupancy logits + scalar predictions
       -> known-value retention + geometry assembly
       -> frozen MetaDiT surrogate
\end{lstlisting}

The original architecture is correctly connected at the data, decoder, and
physics boundaries. Its unresolved scientific issue is target-specific
geometry selection: the requested spectrum must improve decoded geometry rather
than merely alter the latent representation.

\end{document}

